% GENERATED FILE. Edit canonical lessons, then run npm run generate:books.
% canonical-source-hash: fnv1a64:775bf5c16e43b84a
% canonical-lessons: ZH-C139-lu, ZH-C139-huang, ZH-C139-hui, ZH-C139-yuan, ZH-C139-fenhong

\chapter{Colours and Shapes}
\label{ch:zh-colours-and-shapes}
\hlchaptermodality{139}

\begin{chapteropening}
\textbf{By the end of this chapter:} \emph{I can name green, yellow, grey and pink, and say that something is round.}

Complete the last lesson of chapter 139: I can name green, yellow, grey and pink, and say that something is round.
\end{chapteropening}

\section[lǜ]{lǜ — green}

\label{lesson:ZH-C139-lu}



\begin{quote}
Before the new one: say the Mandarin for every single, then the Mandarin for the others.
\end{quote}

\subsection*{You'll want to know: lǜ}
\emph{lǜ} — \textquotedblleft{}green\textquotedblright{}.

\begin{quote}
\emph{chūntiān de cǎo hěn lǜ} — In spring the grass is very green.
\end{quote}

Its characters come later; for now it is a word for the ear.

\begin{cousinweb}
\emph{lǜ} is inherited from Middle Chinese \emph{ljowk} (Guangyun, in Baxter's transcription), glossed 'the colour of blue and yellow'. The regular sound change would give \emph{lù}; why the everyday word is \emph{lǜ} is not settled.
\end{cousinweb}

\begin{grammarlens}[title={What you've built}]
One more word for describing things.
\end{grammarlens}

\subsection*{Your turn}
\begin{itemize}
\raggedright
  \item \emph{Say it:} \emph{lǜ}
  \item \emph{Say it:} \emph{lǜ}, once more
  \item \emph{Say it:} \emph{chūntiān de cǎo hěn lǜ} — In spring the grass is very green.
  \item \emph{Recall:} say the Mandarin for every single, then the Mandarin for the others, then say \emph{lǜ} again
\end{itemize}

\begin{tcolorbox}[breakable,colback=teal!4,colframe=teal!35!black,title={Before you move on}]
What does lǜ mean? (\textquotedblleft{}Green\textquotedblright{}.) Say it once more.
\end{tcolorbox}
\section[huáng]{huáng — yellow}

\label{lesson:ZH-C139-huang}



\begin{quote}
Before the new one: say the Mandarin for the others, then the Mandarin for green.
\end{quote}

\subsection*{You'll want to know: huáng}
\emph{huáng} — \textquotedblleft{}yellow\textquotedblright{}.

\begin{quote}
\emph{zhège píngguǒ shì huáng de} — This apple is a yellow one.
\end{quote}

Its characters come later; for now it is a word for the ear.

\begin{cousinweb}
\emph{huáng} is inherited from Middle Chinese \emph{hwang} (Guangyun, in Baxter's transcription), glossed 'the colour of the centre'.
\end{cousinweb}

\begin{grammarlens}[title={What you've built}]
One more word for describing things.
\end{grammarlens}

\subsection*{Your turn}
\begin{itemize}
\raggedright
  \item \emph{Say it:} \emph{huáng}
  \item \emph{Say it:} \emph{huáng}, once more
  \item \emph{Say it:} \emph{zhège píngguǒ shì huáng de} — This apple is a yellow one.
  \item \emph{Recall:} say the Mandarin for the others, then the Mandarin for green, then say \emph{huáng} again
\end{itemize}

\begin{tcolorbox}[breakable,colback=teal!4,colframe=teal!35!black,title={Before you move on}]
What does huáng mean? (\textquotedblleft{}Yellow\textquotedblright{}.) Say it once more.
\end{tcolorbox}
\section[huī]{huī — grey}

\label{lesson:ZH-C139-hui}



\begin{quote}
Before the new one: say the Mandarin for green, then the Mandarin for yellow.
\end{quote}

\subsection*{You'll want to know: huī}
\emph{huī} — \textquotedblleft{}grey\textquotedblright{}.

\begin{quote}
\emph{wǒ de māo shì huī de} — My cat is a grey one.
\end{quote}

Its characters come later; for now it is a word for the ear.

\begin{cousinweb}
\emph{huī} is inherited from Middle Chinese \emph{xwoj} (Guangyun, in Baxter's transcription), which quotes an older book: 'dead fire', that is, ash. The colour of ash is the likely step to grey, though it is not given there.
\end{cousinweb}

\begin{grammarlens}[title={What you've built}]
One more word for describing things.
\end{grammarlens}

\subsection*{Your turn}
\begin{itemize}
\raggedright
  \item \emph{Say it:} \emph{huī}
  \item \emph{Say it:} \emph{huī}, once more
  \item \emph{Say it:} \emph{wǒ de māo shì huī de} — My cat is a grey one.
  \item \emph{Recall:} say the Mandarin for green, then the Mandarin for yellow, then say \emph{huī} again
\end{itemize}

\begin{tcolorbox}[breakable,colback=teal!4,colframe=teal!35!black,title={Before you move on}]
What does huī mean? (\textquotedblleft{}Grey\textquotedblright{}.) Say it once more.
\end{tcolorbox}
\section[yuán]{yuán — round}

\label{lesson:ZH-C139-yuan}



\begin{quote}
Before the new one: say the Mandarin for yellow, then the Mandarin for grey.
\end{quote}

\subsection*{You'll want to know: yuán}
\emph{yuán} — \textquotedblleft{}round\textquotedblright{}.

\begin{quote}
\emph{zhè zhāng zhuōzi shì yuán de} — This table is a round one.
\end{quote}

Its characters come later; for now it is a word for the ear.

\begin{cousinweb}
\emph{yuán} is inherited from Middle Chinese \emph{hjwen} (Guangyun, in Baxter's transcription), glossed 'the form of heaven'. The plain sense round is not among the glosses, so the step is not settled.
\end{cousinweb}

\begin{grammarlens}[title={What you've built}]
One more word for describing things.
\end{grammarlens}

\subsection*{Your turn}
\begin{itemize}
\raggedright
  \item \emph{Say it:} \emph{yuán}
  \item \emph{Say it:} \emph{yuán}, once more
  \item \emph{Say it:} \emph{zhè zhāng zhuōzi shì yuán de} — This table is a round one.
  \item \emph{Recall:} say the Mandarin for yellow, then the Mandarin for grey, then say \emph{yuán} again
\end{itemize}

\begin{tcolorbox}[breakable,colback=teal!4,colframe=teal!35!black,title={Before you move on}]
What does yuán mean? (\textquotedblleft{}Round\textquotedblright{}.) Say it once more.
\end{tcolorbox}
\section[fěnhóng]{fěnhóng — pink}

\label{lesson:ZH-C139-fenhong}



\begin{quote}
Before the new one: say the Mandarin for grey, then the Mandarin for round.
\end{quote}

\subsection*{You'll want to know: fěnhóng}
\emph{fěnhóng} — \textquotedblleft{}pink\textquotedblright{}.

\begin{quote}
\emph{tā de màozi shì fěnhóng de} — Her hat is a pink one.
\end{quote}

Its characters come later; for now it is a word for the ear.

\begin{cousinweb}
\emph{fěnhóng} is built from \emph{fěn} 'powder' and \emph{hóng} 'red', both inherited from Middle Chinese, \emph{pjunX} and \emph{huwng} (Guangyun, in Baxter's transcription): red as pale as face powder. The dictionary glosses the first with a powder made by burning lead.
\end{cousinweb}

\begin{grammarlens}[title={What you've built}]
One more word for describing things.
\end{grammarlens}

\subsection*{Your turn}
\begin{itemize}
\raggedright
  \item \emph{Say it:} \emph{fěnhóng}
  \item \emph{Say it:} \emph{fěnhóng}, once more
  \item \emph{Say it:} \emph{tā de màozi shì fěnhóng de} — Her hat is a pink one.
  \item \emph{Recall:} say the Mandarin for grey, then the Mandarin for round, then say \emph{fěnhóng} again
\end{itemize}

\begin{tcolorbox}[breakable,colback=teal!4,colframe=teal!35!black,title={Before you move on}]
What does fěnhóng mean? (\textquotedblleft{}Pink\textquotedblright{}.) Say it once more.
\end{tcolorbox}
